\section*{Sources \& Limitations}
\addcontentsline{toc}{section}{Sources \& Limitations}

\subsection*{Primary sources consulted}
9月28日のAnthropicのSonnet 5.5、9月29日のOpenAIのGPT-6.1 Sol、9月30日のGoogleのGemini 4 Argon、9月28日のH CompanyのHolo4、10月2日のELYZAのThinking LLM-jp-4、9月29日のOllamaのv0.35、同日のOpenAIのAgents APIの計算機使いとdotsとDevDayの束ね、9月29日のProvenanceGuardの団子記事、9月28日のNVIDIAのOpen Agent Safety、9月30日のSynthID Bio、9月28日のOpenAIの安全事例集、10月1日のCloudflareのOAuth渡し、9月29日のContext Language Models（arXiv v1とr16の一次注記つき）、10月1日のOlmo-core 3、9月29日のAA-AgentPerf-Local、9月30日のOpen TTS Leaderboard、9月28日のHugging FaceのRL Environments Hub、10月1日のFLUX 3、同日のVSS Blueprint 3.3、9月30日のNemotronのサウジ・アラビア言葉の手直しとNeMo RelayとAMD Ross、9月28日のAMDとWorld Labsの結び、10月1日のCloudflareのAI Search、10月1日のClefとStrands Decider 2Bを範囲として読んだ。本号では全文PDFや手法の細目、図表の詳細までは扱わず、その範囲を超える主張は行わない。

\subsection*{Temporal boundary}
通常の範囲は2026-09-25 18:00から2026-10-02 18:00（いずれもAmerica/New\_York、終了時刻は含まない）。10月1日の集まりの話や10月7日の写しの到着、10月8日までの半額の終わりは範囲の区切りとともに読む。重みの置き場の生まれをもって写しの公開と読まない。日付は日の単位のものが多く、時刻レベルの断定は確かめたものに限る。

\subsection*{Vendor-reported boundary}
測定値や比較、コスト、スループット、利用見通しの説明は、ベンダーや模型の置き場、提携の報告として受け止め、独立した再現とは切り分ける。Sonnet 5.5の70.6\%や1844、Solの値札、Argonの到達点、Holo4やELYZAの表、ClefやStrandsの待ち時間、VSSやASRの催しの計りといった値はベンダー側の報告であり、条件の違いを含めて本号の検証範囲外である。

\subsection*{Community observation boundary}
公開投稿の台帳は原則として関心の方向を示す背景にとどめる。ただし、一次発表として本文で明示的に引用した公式アカウントの投稿は、その発表内容に限って一次資料として扱う。台帳に載っているという事実だけでは、いかなる行も技術的な根拠に格上げしない。独立した投稿やコミュニティの反応は、仕様やベンチマーク、日付、利用条件、価格、ライセンス、アーキテクチャ、公開事実の根拠にはならない。全行の内訳は本号の公開台帳（community-observation-ledger.md）で直接確かめられる。行単位の台帳を正とし、要約文の数えに頼らない。

\subsection*{What was not elevated}
二件の催し記事（AstaBriefとAutoSynthData）は、一次資料の受け入れ整理が終わっていないため今号の章には載せない。いずれも窓内の素材的に重い品であることは認め、載せない理由を隠さない。64GBの品の第三者の入手法は未来の話として、今号の品にはしない。収集の台帳そのものは手筋の覚えに残し、物語には入れない。窓より前の紙一枚は文脈としてのみ扱う。噂の段階の話題は今週の技術的な出来事として取り上げない。
